\section{Introduction}\label{sec:intro}

% Sources of the first sentence (page, sentence of the Introduction):
% - Knill96, p. 1, sentence 2: "The classical machine performs off-line
%   classical computations and controls the evolution of the quantum
%   registers by initializing them in certain preparable states, operating on
%   them with elementary unitary operations and measuring them when needed."
% - FLV25, p. 1, sentences 4 and 5: "Knill's QRAM model [34] describes a
%   generic interface for interacting with such a quantum memory. The memory
%   is modeled with uniquely identified quantum registers, each holding one
%   quantum bit -- also called a qubit."
% - CHPV26, p. 1, sentences 2 and 3 (after the heading "Motivations."): "In
%   this model, a program interacts with a quantum memory through basic
%   operations complying with the laws of quantum mechanics. These operations
%   include a fixed set of quantum gates, chosen to be universal, as well as
%   a probabilistic measurement of qubits."
A quantum program interacts with a memory made of quantum bits,
named qubits~\cite{Knill96,FLV25,CHPV26}. Unlike a
classical bit, a qubit is in a superposition of the values $0$ and $1$.
A program is built from a few operations on qubits: it initialises them, it
transforms them with unitary gates, and it measures them. A measurement
returns a classical outcome, each outcome having some probability, and it
changes the state of the memory.

In a quantum protocol several such programs run concurrently, share a
memory, and exchange messages, typically the outcomes of their measurements.
Process calculi are the standard formal languages for concurrent programs,
and \lqCCS{}~\cite{Ceragioli24popl} is a process calculus for quantum ones: it
extends CCS with values and with operations on qubits. Its meaning is
given by a \emph{reduction semantics}~\cite{SangiorgiWalker01}: a relation
that describes, step by step, how a program and its memory evolve.

A reduction semantics describes computations, not interactions: it does not
say which messages a program is ready to send or to receive, that is, what
it offers to the programs it may be run with. This matters when comparing
programs. In testing theory, a program $Q$ may replace
a program $P$ if $Q$ passes every test that $P$ passes, a test being a
program run in parallel; in the
\emph{must-preorder} of De Nicola and Hennessy~\cite{DNH84}, passing a test
means succeeding in every execution. Owing to the universal quantification
over tests, this definition does not yield a practical proof
method~\cite{BCLS25}. This contextual preorder needs an \emph{alternative
characterisation}: an equivalent relation, defined by looking at the two
programs only. Such characterisations are stated on a \emph{labelled
transition system} (LTS), in which every step carries a label: either the
action that the program offers to its environment, such as sending or
receiving a message, or the silent label $\tau$ for an internal step.

Two recent lines of work call for each other on this point. On the one hand,
testing preorders have been defined for \lqCCS and used to verify a quantum
protocol~\cite{Ceragioli27}; the authors leave their alternative
characterisations as future work, and point to the work of Bernardi et
al.~\cite{BCLS25,BFL26}. On the other hand, Bernardi et al.\ have
characterised the must-preorder for every LTS that satisfies a few axioms,
whatever the calculus, with proofs checked by the Rocq proof
assistant~\cite{BFL26}. The authors see this result as a step towards
quantum process calculi, and expect that the next ones require a more general
notion of duality, which says which actions can meet, and a shared medium
other than a multiset of messages, such as a quantum memory.

The two lines do not meet yet. The existing labelled semantics of
\lqCCS{}~\cite{Ceragioli25,Ceragioli27} is defined on the pairs of a memory
and a program taken as a whole, and one step leads to a probability
distribution of such pairs. In the framework of Bernardi et al.~\cite{BFL26} the target of a transition is
a state, not a distribution of states. Moreover, the parties that interact
there, a program and its test, are two separate LTSs, put together by a
parallel composition of LTSs; treating the memory as a shared medium, as
suggested above, asks for the same separation between a program and its
memory. An LTS of this kind has to be defined for \lqCCS, and it has to be
trusted. The standard criterion is the \emph{Harmony
lemma}~\cite{Milner89,SangiorgiWalker01}: the reductions of a system are
exactly its silent transitions, up to \emph{structural congruence}, the
equivalence that identifies programs differing only in the way they are
written. Its proof is tedious enough to be a challenge of the Concurrent
Calculi Formalisation Benchmark, recently answered in Rocq for a minimal
$\pi$-calculus~\cite{KanLenglet26}. For \lqCCS even its statement is
% Source of the two terms: Faggian22, p. 5:10 (page 10 of the PDF),
% Section 2.4, paragraph "Probabilistic vs Non-deterministic", sentence 1:
% "It is important to understand the distinction between probabilistic choice
% (which globally happens with certitude) and non-deterministic choice (which
% leads to different distributions of outcomes.)"
delicate, because a system makes two kinds of choices, for which we use the
terminology of Faggian~\cite{Faggian22}. A \emph{probabilistic choice} is the outcome of a
measurement: all the outcomes happen, each with its probability. A
\emph{non-deterministic choice} is the selection of the next step among the
ones that are ready: it is made by a scheduler which, to be physically
admissible, must not depend on the state of the
memory~\cite{Ceragioli25}. The question is thus: \emph{is there an LTS for
\lqCCS, built with the parallel composition of LTSs of Bernardi et al.~\cite{BFL26}, whose
silent transitions are exactly the reductions of \lqCCS, with the same
choices of the scheduler and the same probabilities?}

\paragraph{Main contribution.} We answer this question positively. We treat
the memory as one more component of the system: we give an LTS to the memory
$\rho$, another one to the program $P$, which runs without memory, and we
let the LTS of \lqCCS be their parallel composition $\rho \lpar P$. A gate or
a measurement is then an interaction between the two. \Cref{lem:harmony}
proves that
\[
\conf{\rho}{P} \red_{r} \conf{\rho'}{Q}
\iff
\rho \lpar P \step{\tau}{r} \rho' \lpar Q' \text{ for some } Q' \equiv Q ,
\]
where the left-hand side is a reduction, the right-hand side a silent
transition, $\equiv$ the structural congruence, and $r$ the choice of the
scheduler. The result is mechanised in Rocq, on top of the library of
Bernardi et al.~\cite{BFL26}.

\iffalse%
\paragraph{Technical challenges.} To arrive at \cref{lem:harmony} we have to
overcome three issues. The first is to tell the two kinds of choices apart
when a step leads to a single state and not to a distribution. As
in the work of Ceragioli et al.~\cite{Ceragioli25}, every instruction of a program carries a name, its
\emph{tag}, and we let every step carry the choice of the scheduler, in the
style of Chatzikokolakis and Palamidessi~\cite{CP10}: the tag of the instruction that is executed, or the
two tags of a communication. A measurement has one step for each outcome;
all of them carry the same choice, and the probability of each is recorded
in the memory it reaches, as in the semantics of sequential quantum
programs~\cite{UVY26}. The second issue is the composition itself.
In the framework of Bernardi et al.~\cite{BFL26} the two components have the same actions, and a function
maps each action to its dual. A program and a memory do not have the same
actions: we replace this function by a \emph{duality predicate}, a relation
between the actions of the two components. The third issue is that the
Harmony lemma has to preserve the choices, although the structural
congruence may exchange the sender and the receiver of a communication: the
choice of a communication has to be an unordered pair of tags.
\fi

\paragraph{Further contributions and paper outline.} In \cref{sec:prelim} we
recall the quantum notions we need, and we define LTSs with choices and
their composition with respect to a duality predicate, a relation that says
which actions of the two components can meet; neither definition is
specific to quantum computing. \Cref{sec:red} presents \lqCCS, extended with
the initialisation of qubits and with recursion, and its reduction semantics. \Cref{sec:lts}
defines the LTS of programs and the LTS of the memory. The latter has two
kinds of actions only, applying an operation and sending the outcome of a
measurement: the quantum part of the calculus is confined to one small
component.
% , which the mechanisation takes as a parameter.
\Cref{sec:harmony} proves the Harmony lemma. \Cref{sec:conclusion} concludes,
and \cref{sec:future} discusses the next step, the characterisation of a
probabilistic must-preorder for \lqCCS.
